\def\ChapterCode{WUN}
% \begin{savequote}
% \includegraphics[width=3.5cm]{./images/wikipedia-pd/Battlefield_wounds_300px.jpg}
% 
% {\tiny\itshape Hans von Gersdorff, Feldbuch der Wundarznei, 1517.}
% \end{savequote}
\chapter
[\CO{[WUN]} Wunden]
{Wunden}
\label{chp:WUN}
% \AasRsN{RS~VII}

\ChapterIntro
{%
 Eine Wunde ist eine Unterbrechung des
Zusammenhangs von Körpergeweben durch mechanische,
thermische, Strahlen- oder chemische Einwirkungen.
}
{%
\begin{description}
  \item[Maintainer:] Sebastian Gabriel
  \item[Autoren:] Diverse
  \item[Reviewer:] Standard-Reviewprozess
  \item[Version:] {\DocVersion} ({\DocVersionInfo})
  \item[SHA1:] {\DocGitCommit}
\end{description}

{\TxTeilkapitelAASS}
}


\clearpage




\section{Allgemeines und Einteilung von Wunden}

\DefinitionInPara{Eine \T{Wunde} \Z{(\Lang{Lat.} \IX{Vulnus} (\I{Vuln.}))} ist eine Unterbrechung des
Zusammenhangs von Körpergeweben durch mechanische,
thermische, Strahlen- oder chemische Einwirkungen
\cite{Pschy259}.}
\Commentary[Wunde]{Die übliche Beschreibung einer Wunde als eine Verletzung der
Haut ist nicht korrekt. Meist werden mehrere Gewebearten in
Mitleidenschaft gezogen. Definition
\protect\cite{Pschy259}: Unterbrechung
des Zusammenhangs von Körpergeweben mit oder ohne
Substanzverlust.}
\label{topic:wunden-einteilung}%
\label{topic:wunden}%
\index{Wunde}%

Wunden lassen sich \E{nach verschiedenen Gesichtspunkten}
einteilen bzw. unterscheiden. In den folgende Abschnitten
sind Einteilungen nach der physikalischen
Verletzungsursache, nach der Tiefe der Schädigung sowie nach
der Art der Schädigung (Wundart) angeführt.




\TextSlide{Physikalische Verletzungsursache}{%
Eine Einteilung bezieht sich auf die \EE{physikalische
Verletzungsursache} \cite[S.1796]{Pschy259}:

\begin{enumerate}
  \item \T{Mechanische Wunden}:
  \index{Wunden!mechanische}Entstehen durch spitze oder
  stumpfe Gewalteinwirkung auf den Körper (Stichwunde,
  Schürfwunde, Schusswunde, Bisswunde, \ldots)
  \item \T{Thermische Wunden}: \index{Wunden!thermische}
  Entstehen durch Kälte- oder Hitzeeinwirkung auf den Körper
   (Verbrennungen, Stromeinwirkung, Erfrierung, \ldots)
  \item \T{Chemische Wunden}:
  \index{Wunden!chemische}Entstehen durch Verätzungen mit
  Laugen oder Säuren
  \item \T{Strahlenbedingte Wunden}: \index{Wunden!strahlenbedingte} Entstehen
  durch Ultraviolettstrahlung oder ionisierende Strahlung
  (radioaktive Strahlung, Röntgenstrahlung)
\end{enumerate}
}{
\begin{itemize}
  \item Mechanische Wunden
  \item Thermische Wunden
  \item Chemische Wunden
  \item Strahlenbedingte Wunden
\end{itemize}
}{}{}{}



\TextSlide{Tiefe der
Schädigung}{%
Eine weitere Unterteilung kann nach der \EE{Tiefe der
Schädigung} des Gewebes vorgenommen werden.

\begin{itemize}
  \item \EE{Oberflächliche} Wunden: Leichte Verletzungen,
  welche max. bis zur Lederhaut reichen.
  \item \EE{Tiefe} Wunden: Neben der Verletzung der Haut
  können hier auch Muskeln, Sehnen, größere Blutgefäße,
  innere Organe und Knochen betroffen sein \cite{LPN3}.
  \item \EE{Penetrierende} Wunden: Penetrierende Wunden
  sind tiefe Wunden mit Eröffnung mindestens einer
  Körperhöhle (Schädel, Thorax, Bauch) \cite{LPN3}.
  % 	\item \EE{Geschlossene Wunden:} Verletzungen der
     % 	unteren Hautschichten ohne Eröffnung der Oberhaut. (Hämatom) % gestrichen
        % lt. STEUER
\end{itemize}
}{
\begin{itemize}
  \item Oberflächliche Wunden
  \item Tiefe Wunden
  \item Penetrierende Wunden
\end{itemize}
}{}{}{}


\TextSlide{Zeitlicher Verlauf}{%
Eine weitere Unterteilung kann nach dem \EE{zeitlichen Verlauf} 
getroffen werden:

\begin{itemize}
  \item \EE{Frische} Wunden \Z{(Rezente Wunden, \enquote*{rec.})}:
  Wunden innerhalb der ersten 6\Hour* nach der Entstehung
  \item \EE{Nicht frische} Wunden \Z{(Nicht rezente Wunden,
  \enquote*{non rec.})}: Wunden älter als 6\Hour*
  \item \EE{Chronische} Wunden: Wunden, welche seit langer
  Zeit bestehen und Wundheilungsstörungen aufweisen
  % 	\item \EE{Geschlossene Wunden:} Verletzungen der
     % 	unteren Hautschichten ohne Eröffnung der Oberhaut. (Hämatom) % gestrichen
        % lt. STEUER
\end{itemize}
}{
\begin{itemize}
  \item Frische Wunden: $\leq$\,6\Hour*
  \item Nicht-frische Wunden:
  $>$\,6\Hour*
  \item Chronische Wunden
\end{itemize}
}{}{}{}





\TextSlide{Wundarten}
{%
Unterschiedliche \E{Entstehungsmechanismen} erzeugen
unterschiedliche \E{Wundarten} bzw. \E{Wundformen}.
In \FullPrettyRef{tab:Wundarten} findet sich eine Aufstellung 
und kurze Erklärung der häufigsten Wundarten.

}{
\FullPrettyRef{tab:Wundarten}
}
{}
{}
{}

\Layout{57}{}
{%

}{
}{\index{Wundarten}\label{tab:Wundarten}
\begin{tabulary}{\linewidth}{lLL}\addlinespace\toprule\addlinespace
% \addlinespace
\noindent\TabHeading{Wundart} & \noindent\TabHeading{Beschreibung} & \noindent\TabHeading{\Z{Lat. Diagnose}}
\\\addlinespace\midrule\addlinespace
\noindent\T{Schürfwunde} 
&
Oberflächliche Wunde
& 
\noindent\Z{\IX{Excoriatio} (\I{Excor.})}
\\\addlinespace
\noindent\T{Schnittwunde}
&
Glatte Wundränder; tiefer liegende Nerven und Blutgefäße
können verletzt sein
& 
\noindent\Z{\IX[Vulnus!scissum]{Vulnus scissum} (\I[Vuln.!sciss.]{Vuln. sciss.})}
\\\addlinespace
\noindent\T{Stichwunde}
&
Glatte Wundränder; meist kleine Wunde dafür aber tiefer,
kann mitunter in Körperhöhlen penetrieren; Das harmlose
Aussehen kann über das tatsächliche Verletzungsausmaß hinweg
täuschen!
& 
\noindent\Z{\IX[Vulnus!ictum]{Vulnus ictum} (\I[Vuln.!ict.]{Vuln. ict.})}
\\\addlinespace 
\noindent\T{Risswunde}
&
Unregelmäßige Wundränder; meist nur Schädigung der Haut und
	keine tieferliegenden Nerven und Blutgefäßen
& 
\noindent\Z{\IX[Vulnus!laceratum]{Vulnus laceratum} (\I[Vuln.!lac.]{Vuln. lac.})}
\\\addlinespace
\noindent\T{Quetschwunde}
&
Unregelmäßige Wundränder; es kann zum Aufplatzen der Haut kommen (Platzwunde)
& 
\noindent\Z{\IX[Vulnus!contusum]{Vulnus contusum} (\I[Vuln.!cont.]{Vuln. cont.})}
\\\addlinespace
\noindent\T{Rissquetschwunde} (\I{RQW})
&
Kombination aus Riss- und Quetschwunde
& 
\noindent\Z{\IX[Vulnus!lacero-contusum]{Vulnus lacero-contusum} (\I{Vlc.})}
\\\addlinespace
\noindent\T{Bisswunde} 
&
Kombination aus Riss-, Stich- und Quetschwunde mit hoher Infektionsgefahr
& 
\noindent\Z{\IX[Vulnus!morsum]{Vulnus morsum (animalis, hominis)} (\I[Vuln.!mors.]{Vuln. mors. (anim., hom.)})}
\\\addlinespace
\noindent\T{Ablederung} 
&
Abscherung und Ablederung der Haut infolge von Rotation und Quetschung, Ablösung der Haut vom Untergrund
& 
\noindent\Z{\IX[Decollement]{Decollement}}
\\\addlinespace
\noindent\T{Amputation} 
&
Abtrennung eines Körperteiles
& 
\noindent\Z{\IX{Amputatio} (\I{Amp.})}
\\\addlinespace
\noindent\T{Pfählung} 
&
Eindringen (oder Durchdringen) und Verbleiben von spitzen,
u.\,U. unregelmäßigen Fremdkörpern (z.\,B. Messer,
Injektionsnadel) 
& 
\noindent\Z{\IX[Vulnus!perforans]{Vulnus perforans}}
\\\addlinespace 
\noindent\T{Schusswunde} 
&
Meist klein aussehende Wunde, das eigentliche Schadensausmaß im Körperinneren
ist meist nicht ersichtlich (Druckwelle kann starke innere
Verletzungen verursachen),
\Z{evtl. Ausschuss größer als Einschuss}
& 
\noindent\Z{\IX[Vulnus!sclopetarium]{Vulnus sclopetarium} (\I[Vuln!sclopet.]{Vuln. sclopet.})}
\\\addlinespace
\noindent\T{Verbrennung} 
&
Durch Hitzeeinwirkung entstandene Schädigung, siehe \FullPrettyRef{topic:verbrennung}
& 
\noindent\Z{\IX{Combustio} (\I{Comb.})}
\\\addlinespace
\noindent\T{Verbrühung} 
&
Sonderform der Verbrennung, verursacht durch heiße Flüssigkeiten oder Dampf
& 
\noindent\Z{\IX{Ambustio} (\I{Amb.})}
\\\addlinespace
\noindent\T{Verätzung} 
&
Lokale Schädigung der
Oberfläche durch stoffe aufgrund einer irreversiblen Zerstörung (Denaturierung) von
Eiweißstoffen, siehe \FullPrettyRef{topic:veraetzung-intoxikation}
& 
\noindent\Z{\IX{Cauterisatio} (\I{Cauteris.})}
\\\addlinespace
\noindent\EE{(Hämatom)} 
&
Bluterguss; Ansammlung von Blut im Gewebe außerhalb der Blutgefäße 
& 
\noindent\Z{\IX{Haematoma} (\I{Haem.})}
\\\addlinespace\bottomrule
\end{tabulary}
}{Wundarten.}{}

% \begin{PictureSeriesWide}{}{Bilderserie: Wunden: \EE{Verätzungen und
% Verbrennungen}}
% \PictureSeriesRowWide{3}
% {./images/trauma/burn-chem01.jpg}
% {Verätzung
% \SourcePicture{Hauer}}
% {./images/trauma/burn-chem02.jpg}
% {Verätzung
% \SourcePicture{Hauer}}
% {./images/trauma/combustio-sa-02-hand.jpg}
% {Ablederung nach schwerer
% Verbrennung
% \SourcePicture{Hauer}}
% {}
% {}
% \end{PictureSeriesWide}

\begin{PictureSeriesWide}{}{Bilderserie: \SlideHeading{Wunden}}
\PictureSeriesRowWide{3}
{./images/wunde_rqw1.jpg}
{Rissquetschwunde
vor der Versorgung im Spital
\SourcePicture{Hauer}}
{./images/wunde_rqw2.jpg}
{Rissquetschwunde
nach der Versorgung im Spital
\SourcePicture{Hauer}}
{./images/pulsader_schnitt.jpg}
{Schnittverletzung
oberhalb der \enquote{Pulsader}
\SourcePicture{Hauer}}
{}
{}
\PictureSeriesRowWide{3}
{./images/trauma/stichwunde-klein.jpg}
{Glatte
Wundränder und eigentlich ganz unauffällig: Die Stichwunde.
\SourcePicture{Hauer}}
{./images/pneumothorax_stichverletzung_diskret.jpg}
{Stichverletzung:
unauffällig. \Ca{Dieser Patient ist
lebensgefährlich verletzt!}
(Pneumothorax)\SourcePicture{Hauer}}
{./images/trauma/stichwunde-thorax-tief-01}
{Stichwunde
mit eröffneter Brusthöhle.
\SourcePicture{Hauer}}
{}
{}
\PictureSeriesRowWide{3}
{./images/trauma/burn-chem01.jpg}
{Verätzung
\SourcePicture{Hauer}}
{./images/trauma/burn-chem02.jpg}
{Verätzung
\SourcePicture{Hauer}}
{./images/trauma/combustio-sa-02-hand.jpg}
{Ablederung nach schwerer
Verbrennung
\SourcePicture{Hauer}}
{}
{}
\end{PictureSeriesWide}



\section{Gefahren bei Wunden}

\TextSlide{\TxGefahren}{
Jede Wunde ist nicht nur eine lokale Verletzung sondern hat
auch (in unterschiedlichem Ausmaß) eine Auswirkung auf den
Gesamtorganismus. Neben der Beurteilung der Wunde selbst
muss daher immer auch eine Beurteilung der Vitalfunktionen
und der Schmerzen des Patienten sowie der
Infektionsgefahren bei der Wundversorgung erfolgen.
\begin{itemize}
  \item \EE{Verletzung wichtiger Gewebestrukturen (Organe):}
  Je nach betroffener Körperregion treten unterschiedliche
  Verletzungsmuster auf.  Daraus resultieren spezielle
  traumatologische Notfallbilder (z.\,B.
  Schädel-Hirn-Trauma, Thorax-Trauma, Bauchtrauma etc.), die in den folgenden
  Abschnitten beschrieben werden. 
  \item \EE{Starke Blutung:} Bei großflächigen Wunden oder
  Verletzung wichtiger Blutgefäße kann es zu starken
  Blutungen und infolge dessen zu einem Volumenmangelschock
  kommen. Blutstillung und Schockbekämpfung sind  die
  wichtigsten Maßnahmen. 
  \item \EE{Starke Schmerzen:} Treten
  im Zusammenhang mit einer Verletzung starke Schmerzen auf,
  muss ein Notarzt für eine Schmerztherapie gerufen werden.
  \item \EE{Eintritt von Krankheitserregern
  (Infektionsgefahr):} Jede Wunde ist eine Eintrittsstelle
  für Krankheitserreger in unseren Organismus. Der Patient
  sollte daher über einen aufrechten \EE{Tetanus}-Impfschutz
  verfügen (s.
  \prettyref{topic:tetanus-impfschutz})\footnote{Zur
  Verifizierung eines aufrechten Tetanus-Impfschutzes muss
  jede Wunde einem Arzt vorgestellt werden.  Eine Impfung
  muss spätestens alle 10 Jahre aufgefrischt werden,  bei
  entsprechenden Verletzungen schon
  früher. \cite{BoehmerTaschRett6}}.  Bei Bisswunden kann
  über den Speichel einer erkrankten Tieres \EE{Tollwut}
  übertragen werden. Bei Bisswunden soll daher versucht
  werden (am Einsatzprotokoll) zu dokumentieren, wer der
  Besitzer des Tieres ist. Unter Umständen kann es notwendig
  sein die Polizei
  einzuschalten.
  %\item EMHO: Kompartmentsyndrom hier sinnvoll?
\end{itemize}
}{
\begin{itemize}
  \item Verletzung wichtiger Organe
  \item Starke Blutung
  \item Starke Schmerzen
  \item Infektion
\end{itemize}
}{}{}{}





\section{Wundversorgung}
Die Wundversorgung orientiert sich an der Wundart sowie an
der Stärke der Blutung. Die Techniken zur Blutstillung
(Hochhalten der betroffenen Gliedmaße, Zudrücken, Abdrücken
der blutzuführenden Arterie, Druckverband und Abbinden) sind
unter \FullPrettyRef{chp:EHI}, beschrieben. Der
professionelle Helfer muss neben den Erste-Hilfe-Maßnahmen
zusätzlich Prioritäten bzgl.
der optimalen Vorgehensweise setzen, die vitale Bedrohung
durch Flüssigkeitsmangel (durch Blutverlust oder
Verbrennungen) einschätzen, die Infektionsgefahren so gut
wie möglich abwenden und ggf. einen Notarzt für eine
Schmerztherapie nachfordern.



% M %%%%%%%%%%%%%%%%%%%%%%%%%%%%%%%%%%%%%%%%%%%%%%%%%%% M %
% Wundversorgung
\ProcedurePrintDefault{TY61000C}




% \begin{itemize}
%   \item Wunde nicht direkt berühren (außer für rasche
%   Blutstillung)
%   \item Steril arbeiten
%   \item Zurückhaltende Wundreinigung (Abspülen mit
%   \ce{NaCl}-Lösung)
%   \item Wundreinigung von innen nach außen
%   \item Fremdkörper in Wunde belassen
%   \item Ausgetretene Strukturen belassen und feucht halten
%   \item Wundverband
% \end{itemize}


\begin{figure}[H]
\begin{center}
\includegraphics[width=0.5\linewidth]{./images/hirtler-aass/frische-Wunde-desinfizieren.pdf}

\Captionof{figure}{Reinigung einer frischen Wunde. Wenn
erforderlich, muss \EE{von der Wunde weg} gereinigt werden,
um keine Keime in die Wunde zu bringen. \SourcePicture{Hirtler.}}
\end{center}
\end{figure}





\section{Dekubitus und Dekubitusprophylaxe}
\index{Dekubitus!Prophylaxe}
\label{topic:dekubitusprophylaxe}



\TextSlide
{\TxBeschreibung}
{%
\DefinitionInPara{\T{Dekubitus}: \I{Wundliegen}. Bereich
lokalisierter Schädigung der Haut und des darunterliegenden Gewebes als Folge von 
lang anhaltendem Druck sowie Scher- und Reibungskräften.}
Durch den Auflagedruck des Körpers kann die 
Durchblutung der Haut abgedrückt werden. Dies geschieht dann,
wenn der Druck, welcher auf der Körperstelle lastet, den
Blutdruck in den Kapillargefäßen übersteigt (ca. 30\MmHg*).
Es kommt dann zur Unterversorgung des Gewebes mit Sauerstoff
und anderen Nährstoffen. Ein gesunder Mensch kann effektiv 
die Entstehung eines Dekubitus verhindern, indem er
sich instinktiv umlagert, sobald man einen Schmerz infolge
einer ungünstigen Lagerung verspürt. Bei bettlägrigen oder 
sogar bewusstlosen oder narkostisierten
Patienten ist dies nur mehr eingeschränkt oder gar nicht
mehr möglich, es obliegt dann dem Fachpersonal für die
korrekte Lagerung zu sorgen.

Bereits der Auflagedruck durch das
Gewicht des eigenen Körpers auf einer Unterlage kann
ausreichen, um zu solch einer
Unterversorgung zu führen. Besonders Stellen, an denen
\EE{Knochen ohne Muskel- oder Fettpolster direkt unter der
Haut} liegen sind gefährdet:
\begin{itemize}
  \item Kreuzbein
  \item Fersen
  \item Sitzbein
  \item Großer Rollhügel (Trochanter maior des
  Oberschenkels)
\end{itemize}
Gleiches gilt für alle Arten
von \EE{Fremdkörpern}, auf denen der Patient aufliegt,
z.\,B.
Kleidungsfalten, Schnallen, im Bett vergessenes
Material, \ldots 

}
{\EE{Dekubitus}: \E{Wundliegen}. Bereich
lokalisierter Schädigung der Haut und des darunterliegenden Gewebes als Folge von 
lang anhaltendem Druck sowie Scher- und Reibungskräften.
\begin{itemize}
  \item Auflagedruck!
  \item Besonders gefährdet:
  \begin{itemize}
    \item Kreuzbein
    \item Fersen
    \item Sitzbein
    \item Großer Rollhügel (Trochanter maior des
    Oberschenkels)
  \end{itemize}
\end{itemize}
}
{}
{}
{}


\TextSlide
{Faktoren}
{Die Entstehung eines Dekubitus wird im wesentlichen von fünf
Faktoren bestimmt:


\begin{itemize}
  \item \EE{Druck}: Je höher der Druck, desto größer ist die
  Schädigung.
  \item \EE{Zeit}: Je länger der Druck auf das Gewebe
  einwirkt, desto größer ist die Schädigung. Erste
  Anzeichen eines Dekubitus können sich bereits nach
  weniger als einer Stunde bilden!
  \item \EE{Scherkraft} und \EE{Reibung}: Scherkräfte bewirken eine
  Verschiebung zwischen den Gewebeschichten und können
  ebenso zu ernsthaften Schädigungen führen. Reibung, z.\,B.
  durch Ziehen des Patienten über das Bett oder einen Sessel, kann Verletzungen der
  oberen Hautschichten verursachen.
  \item \EE{Gewebetoleranz}: Das Gewebe kann unterschiedlich
  anfällig für Verletzungen und Druckgeschwüre sein. Dies
  ist z.\,B. abhängig vom Alter des Patienten,
  Begleiterkrankungen, dem Grad der Bettlägrigkeit,
  Lähmungen und den eingenommenen Medikamenten.
  \item \EE{Mobilität und Empfindungsvermögen}: Je
  eingeschränkter die Mobilität bzw. die (Schmerz-) Empfindungen des Patienten
  sind, desto größer ist das Risiko für einen Dekubitus.
\end{itemize}}
{\begin{itemize}
  \item {Druck}
  \item {Zeit}
  \item {Scherkraft} und {Reibung}
  \item {Gewebetoleranz}
  \item {Mobilität und Empfindungsvermögen}
\end{itemize}}
{}
{}
{}


\Text
{Grade}
{\Z{Einteilung der Schweregrade \cite{Shea1975}:
\begin{itemize}
  \item \E{Grad 1}: Nicht wegdrückbare, umschriebene
  \EE{Hautrötung} bei intakter Haut. Weitere klinische
  Zeichen können \EE{Ödembildung, Verhärtung und eine lokale
  Überwärmung} sein.
  \item \E{Grad 2}: Teilverlust der Haut; Epidermis bis hin
  zu Anteilen der Lederhaut sind geschädigt. Der
  Druckschaden ist oberflächlich und kann sich klinisch als
  \EE{Blase, Hautabschürfung oder flaches Geschwür}
  darstellen.

  \item \E{Grad 3}: Verlust aller Hautschichten
  einschließlich Schädigung oder Nekrose des subkutanen
  Gewebes, die bis auf, aber nicht unter, die
  darunterliegende Faszie reichen kann. Der Dekubitus zeigt
  sich klinisch als \EE{tiefes, offenes Geschwür}.

  \item \E{Grad 4}: \EE{Verlust aller Hautschichten} mit
  ausgedehnter Zerstörung, Gewebsnekrose oder Schädigung von
  \EE{Muskeln, Knochen oder stützenden Strukturen} wie
  Sehnen oder Gelenkkapseln, mit oder ohne Verlust aller
  Hautschichten.

\end{itemize}}}
{}
{}
{}
{}

\TextSlide
{Lagerung bei Dekubitus und zur Prophylaxe}
{
\begin{itemize}
  \item \T{30\textdegree{}-Lagerung}: Der Patient liegt auf
  einem oder zwei weichen Kissen, die unter der Körperhälfte
  eingebracht werden, der Kopf ist durch ein kleines Kissen
  gestützt. Bei korrekter Durchführung kann man leicht unter
  das Kreuzbein oder den Trochanter des Oberschenkelknochens
  fassen. \cite{ThiemesPflege12}
  \item \T{135\textdegree{}-Lagerung}: 
  Von der Seitenlage ausgehend wird der Patient auf eine
  zusammengerollte Decke oder zwei Pölster gelagert. Die
  Lagerung ist korrekt, wenn der unten liegende Trochanter
  ohne Druckbelastung ist. Der unten liegende Fuß kann
  durch einen schmalen Polster freigelagert werden.
  \cite{WundversorgungPflege2}
  
  \item \T{5-Kissen-Weichlagerung}: Der Patient wird komplett
  auf 5 oder mehr Pölstern gelagert. Ziel ist das Freilegen
  bereits geschädigter Bereiche. \cite{ThiemesPflege12}
  \item \II{Bauchlage}
\end{itemize}
Abhängig von Begleiterkrankungen oder -verletzungen können
u.\,U. nicht alle Lagerungsarten bei einem Patienten
angewendet werden.
}
{
\begin{itemize}
  \item {30\textdegree{}-Lagerung}
  \item {135\textdegree{}-Lagerung}
  \item {5-Kissen-Weichlagerung}
  \item {Bauchlage}
\end{itemize}

 }
{}
{}
{}

\Text
{Dekubitusprophylaxe allgemein}
{
\begin{itemize}
  \item \EE{Mobilisation} \ldots{} ist die Maßnahme der
  Wahl. Sie hat zum Ziel, den Menschen in Bewegung zu
  versetzen. Das geschieht nicht durch passiven Transfer oder Umlagerung.
  Nur wenn die Anregungen zu eigenen Bewegungen (wie minimal
  auch immer) erfolglos bleiben, sind passive
  Positionsveränderungen sinnvoll. Bezüglich des
  Zeitintervalls gibt es keine allgemein gültigen Aussagen.
  Sinnvoll ist es regelmäßig zu überprüfen, ob Zeichen von Druckeinwirkungen auftreten
  und das Intervall entsprechend anzupassen.
  \item \EE{Druckentlastung}: Es gibt verschiedene
  Möglichkeiten, die Druckbelastung einzelner Hautregionen zu mindern. Die
  Aktivierung und Mobilisation der Patienten sollte die
  erste Maßnahme sein. Hilfsmittel zur Lagerung und
  Druckentlastung können die Mobilisierung ergänzen. Die
  Maßnahmen zur Druckentlastung sind unabhängig vom Grad der
  Dekubitusausprägung immer durchzuführen, nach Möglichkeit
  die absolute Druckentlastung der Wunde durch Hohl- bzw.
  Freilagerung, beispielsweise mittels Bauchlagerung, aber
  auch durch regelmäßigen Lagewechsel und Vergrößerung der
  Auflagefläche durch spezielle Hilfsmittel wie spezielle
  Wechseldruck-Matratzen oder Luftstrom-Matratzensysteme.
  \cite{Pflegewiki:Dekubitus}
  \item \Z{\EE{Ernährung}: Hierbei sei auf die
  entsprechende Pflege-Fachliteratur verwiesen.}
\end{itemize}


}
{}
{}
{}
{}




\Text
{Dekubitusprophylaxe im Rettungs- und
Krankentransportdienst}
{
Grundsätzlich gilt das vorher Gesagte auch im Rettungs- und
Krankentransportdienst, jedoch sind die technischen
Möglichkeiten eingeschränkt. Bei jedem Krankentransport,
welcher \EE{länger als eine Stunde} dauert, muss jedenfalls
an die Möglichkeit eines Dekubitus gedacht werden und ggfs.
entsprechende Maßnahmen ergriffen werden. Im Rahmen der
Übergabe des Patienten müssen etwaige Dekubitus-Gefahren
oder bereits vorhandene Geschwüre besprochen werden.

Die Mobilisation bzw. regelmäßige Umlagerung soll nach
Möglichkeit auch während des Transports durchgeführt
werden. Als einfaches technisches Hilfsmittel steht die
(nicht abgesaugte!) Vakuummatratze zur Verfügung, welche bei
geringen Dekubitusrisiko als Unterlage verwendet werden
kann. Außerdem sind diverse künstliche
Auflagefelle\ZFootnote{Felle als Unterlage sind im
stationären Bereich eher verpönt, im Rahmen des
Krankentransports haben sie dennoch eine gewisse
Berechtigung.} am Markt verfügbar. Bei Hochrisikopatienten
bzw. ausgedehnten Transportstrecken soll unbedingt mit einer
diplomierten Pflegeperson Rücksprache gehalten werden und
eine individuelle Lösung gefunden werden.

}
{}
{}
{}
{}

\begin{PictureSeriesWide}{}{Bilderserie: \EEE{Dekubitus}}
\PictureSeriesRowWide{2}
{./images/wikipedia-pd/Escarre_Stade_4_detail}
{Dekubitus Grad 4 \SourcePicture{WMC/PD}}
{empty}
{}
{empty}
{}
{empty}
{}
\end{PictureSeriesWide}

\begin{Literary}
\EE{Exzellentes, umfassendes Werk über die Wundbehandlung,
nicht nur für die Pflege:}
\fullcite{WundversorgungPflege2}

\fullcite{ThiemesPflege12}

\fullcite{AerztlicheWundversorgung2}



\EE{Historische Abhandlung:}
\fullcite{GeschichteNeuerenDeutschenChirurgie}

\minisec{Weiters}

\cite{SiewertChirurgie8}

\end{Literary}


\nocite{ThiemesPflege12,WundversorgungPflege2}

\ChapterBibliography

\ChapterEnd